% ============================================================
%  EXAMEN TEMA 6: FUERZAS, ENERGÍA, TABLA PERIÓDICA Y REACCIONES
%  Curso Introductorio — 3er Año
%  Duración: 90 minutos
%  Temas: Leyes de Newton, termodinámica, tabla periódica,
%         modelo atómico y ecuaciones químicas
%  Autor: Ing. Gabriel Astudillo
%  Nota: enunciado limpio (sin pistas ni desarrollo). Las
%  soluciones están en la "Clave de Respuestas" al final.
% ============================================================

\documentclass[12pt, letterpaper]{article}

% ── Paquetes ─────────────────────────────────────────────────

\usepackage{mathwizards-guia}
\mwguia{Examen — Tema 6: Fuerzas, Termodinámica y Química}{Ing. Gabriel Astudillo $\cdot$ Curso 3er Año}

\begin{document}
% ============================================================

% ── Portada / Título ─────────────────────────────────────────
\mwportada{Examen del Tema 6}{Fuerzas, Energía, Tabla Periódica y Reacciones}{Leyes de Newton, Termodinámica, Átomo y Ecuaciones Químicas}

\vspace{4pt}
\begin{cuadroinfo}
  \small
  \textbf{Nombre:} \rule{0.55\linewidth}{0.4pt} \quad \textbf{Fecha:} \rule{0.15\linewidth}{0.4pt} \\[6pt]
  \textbf{Tiempo disponible:} 90 minutos \quad$\bullet$\quad \textbf{Puntaje total:} 100 puntos \\[4pt]
  \textbf{Instrucciones:}
  \begin{enumerate}[leftmargin=*]
    \item Lee cada ejercicio con calma antes de resolverlo.
    \item Muestra \emph{todos} tus pasos, los diagramas de cuerpo libre y las unidades.
    \item Justifica siempre tus clasificaciones y tus respuestas conceptuales.
    \item Usa $g = 9{,}8\;\text{m/s}^2$ salvo que se indique otra cosa.
    \item No se permite el uso de calculadora.
    \item Escribe con letra clara y revisa tu examen antes de entregarlo.
  \end{enumerate}
\end{cuadroinfo}

\vspace{8pt}

% ============================================================
\section{Leyes de Newton (28 puntos)}
% ============================================================

\begin{enumerate}[leftmargin=*]
  \item \textbf{(6 pts)} Asocia cada situación con la ley de Newton correspondiente (1ª, 2ª o 3ª) y justifica brevemente:
        \begin{enumerate}[label=\alph*)]
          \item Un pasajero se inclina hacia adelante cuando el autobús frena bruscamente.
          \item Al empujar una pared, sientes que la pared te empuja a ti.
          \item Un mismo empujón acelera más a un carrito vacío que a uno cargado.
        \end{enumerate}

  \item \textbf{(8 pts)} Sobre un objeto de $4\;\text{kg}$ que se desliza por una superficie horizontal actúan una fuerza aplicada de $30\;\text{N}$ hacia la derecha y una fuerza de fricción de $10\;\text{N}$ hacia la izquierda.
        \begin{enumerate}[label=\alph*)]
          \item Dibuja el diagrama de cuerpo libre.
          \item Calcula la fuerza neta.
          \item Calcula la aceleración del objeto.
        \end{enumerate}

  \item \textbf{(7 pts)} Se desea que un objeto de $6\;\text{kg}$ acelere a $4\;\text{m/s}^2$ sobre una superficie con fricción de $5\;\text{N}$.
        \begin{enumerate}[label=\alph*)]
          \item \textquestiondown Cuál debe ser la fuerza neta?
          \item \textquestiondown Qué fuerza aplicada se necesita?
        \end{enumerate}

  \item \textbf{(7 pts)} Una persona de $70\;\text{kg}$ está dentro de un ascensor. Calcula la fuerza normal del piso cuando el ascensor:
        \begin{enumerate}[label=\alph*)]
          \item sube con aceleración de $1{,}5\;\text{m/s}^2$;
          \item baja con aceleración de $1{,}5\;\text{m/s}^2$.
        \end{enumerate}
\end{enumerate}

\newpage

% ============================================================
\section{Termodinámica Básica (24 puntos)}
% ============================================================

\begin{enumerate}[leftmargin=*]
  \item \textbf{(6 pts)}
        \begin{enumerate}[label=\alph*)]
          \item Explica la diferencia entre \textbf{calor} y \textbf{temperatura}.
          \item \textquestiondown En qué dirección fluye el calor espontáneamente?
          \item Nombra el cambio de estado: sólido a líquido, líquido a gas y sólido a gas.
        \end{enumerate}

  \item \textbf{(8 pts)} Calcula el calor necesario para calentar $300\;\text{g}$ de agua de $25\;^\circ\text{C}$ a $75\;^\circ\text{C}$. Usa $c_{\text{agua}} = 1\;\text{cal/(g}\cdot^\circ\text{C)}$.

  \item \textbf{(10 pts)} Un trozo de hierro de $200\;\text{g}$ a $100\;^\circ\text{C}$ se sumerge en $400\;\text{g}$ de agua a $20\;^\circ\text{C}$. Calcula la temperatura final de equilibrio. Usa $c_{\text{hierro}} = 0{,}11\;\text{cal/(g}\cdot^\circ\text{C)}$ y $c_{\text{agua}} = 1\;\text{cal/(g}\cdot^\circ\text{C)}$.
\end{enumerate}

\newpage

% ============================================================
\section{La Tabla Periódica y el Modelo Atómico (24 puntos)}
% ============================================================

\begin{enumerate}[leftmargin=*]
  \item \textbf{(8 pts)} Un átomo de sodio (Na) tiene número atómico $Z = 11$ y número de masa $A = 23$.
        \begin{enumerate}[label=\alph*)]
          \item \textquestiondown Cuántos protones, neutrones y electrones tiene el átomo neutro?
          \item \textquestiondown El sodio es metal, no metal o metaloide?
          \item El carbono-14 tiene $Z = 6$. \textquestiondown Cuántos neutrones tiene?
        \end{enumerate}

  \item \textbf{(8 pts)} Clasifica cada elemento como metal, no metal o metaloide:
        \[
          \text{hierro}, \quad \text{oxígeno}, \quad \text{silicio}, \quad \text{oro}, \quad \text{carbono}, \quad \text{azufre}
        \]

  \item \textbf{(8 pts)}
        \begin{enumerate}[label=\alph*)]
          \item Ordena cronológicamente los modelos atómicos de Dalton, Thomson, Rutherford y Bohr.
          \item \textquestiondown Qué descubrió Rutherford con su experimento de la lámina de oro?
        \end{enumerate}
\end{enumerate}

\newpage

% ============================================================
\section{Transformaciones y Ecuaciones Químicas (24 puntos)}
% ============================================================

\begin{enumerate}[leftmargin=*]
  \item \textbf{(8 pts)} Clasifica cada transformación como física (F) o química (Q):
        \[
          a)\ \text{hervir agua} \quad b)\ \text{quemar papel} \quad c)\ \text{romper un vidrio}
        \]
        \[
          d)\ \text{oxidar un clavo} \quad e)\ \text{disolver azúcar} \quad f)\ \text{cocinar un huevo}
        \]

  \item \textbf{(8 pts)} Para la ecuación de la combustión del metano:
        \[
          \text{CH}_4\text{(g)} + 2\,\text{O}_2\text{(g)} \rightarrow \text{CO}_2\text{(g)} + 2\,\text{H}_2\text{O(g)}
        \]
        \begin{enumerate}[label=\alph*)]
          \item Identifica los reactivos y los productos.
          \item Indica el estado físico de cada sustancia.
          \item Describe la ecuación con palabras.
        \end{enumerate}

  \item \textbf{(8 pts)}
        \begin{enumerate}[label=\alph*)]
          \item Escribe la ecuación química de la electrólisis del agua (agua líquida produce hidrógeno y oxígeno gaseosos), incluyendo los estados físicos.
          \item Menciona \textbf{tres} indicios que permiten sospechar que ha ocurrido una transformación química.
        \end{enumerate}
\end{enumerate}

\newpage

% ============================================================
\ifmwclaves
  \section{Clave de Respuestas}
  % ============================================================

  \subsection*{Sección 1: Leyes de Newton}

  \begin{enumerate}[leftmargin=*, label=\textbf{\arabic*.}]
    \item a) \textbf{1ª ley (inercia)}: el pasajero tiende a mantener su movimiento hacia adelante mientras el autobús frena.
          \quad b) \textbf{3ª ley (acción-reacción)}: la fuerza que ejerces sobre la pared se devuelve con igual magnitud y sentido opuesto.
          \quad c) \textbf{2ª ley ($F = ma$)}: con la misma fuerza, mayor masa produce menor aceleración.

    \item a) DCL: peso ($\vec{W}$) hacia abajo, normal ($\vec{N}$) hacia arriba, fuerza aplicada ($\vec{F}$) hacia la derecha, fricción ($\vec{f}$) hacia la izquierda.
          \quad b) $F_{\text{neta}} = 30 - 10 = 20\;\text{N}$ hacia la derecha.
          \quad c) $a = \dfrac{F_{\text{neta}}}{m} = \dfrac{20}{4} = 5\;\text{m/s}^2$.

    \item a) $F_{\text{neta}} = m\,a = 6 \times 4 = 24\;\text{N}$.
          \quad b) $F_{\text{aplicada}} = F_{\text{neta}} + f = 24 + 5 = 29\;\text{N}$.

    \item a) Ascensor subiendo con aceleración: $N = m(g + a) = 70(9{,}8 + 1{,}5) = 70(11{,}3) = 791\;\text{N}$.
          \quad b) Ascensor bajando con aceleración: $N = m(g - a) = 70(9{,}8 - 1{,}5) = 70(8{,}3) = 581\;\text{N}$.
  \end{enumerate}

  \newpage

  \subsection*{Sección 2: Termodinámica Básica}

  \begin{enumerate}[leftmargin=*, label=\textbf{\arabic*.}]
    \item a) La \textbf{temperatura} mide la agitación de las partículas de un cuerpo; el \textbf{calor} es la energía que se transfiere entre cuerpos a distinta temperatura.
          \quad b) El calor fluye espontáneamente del cuerpo \textbf{caliente al frío}.
          \quad c) Sólido a líquido: \textbf{fusión}. Líquido a gas: \textbf{vaporización}. Sólido a gas: \textbf{sublimación}.

    \item $Q = m\,c\,\Delta T = 300 \times 1 \times (75 - 25) = 300 \times 50 = 15\,000\;\text{cal}$.

    \item En el equilibrio: $Q_{\text{cedido por el hierro}} = Q_{\text{absorbido por el agua}}$.
          \[
            200(0{,}11)(100 - T_f) = 400(1)(T_f - 20) \Rightarrow 22(100 - T_f) = 400(T_f - 20).
          \]
          \[
            2\,200 - 22T_f = 400T_f - 8\,000 \Rightarrow 10\,200 = 422T_f \Rightarrow T_f \approx 24{,}2\;^\circ\text{C}.
          \]
  \end{enumerate}

  \newpage

  \subsection*{Sección 3: La Tabla Periódica y el Modelo Atómico}

  \begin{enumerate}[leftmargin=*, label=\textbf{\arabic*.}]
    \item a) Protones $= 11$, neutrones $= A - Z = 23 - 11 = 12$, electrones $= 11$ (átomo neutro).
          \quad b) El sodio es un \textbf{metal}.
          \quad c) Carbono-14: neutrones $= 14 - 6 = 8$.

    \item Hierro: metal. \quad Oxígeno: no metal. \quad Silicio: metaloide. \quad Oro: metal. \quad Carbono: no metal. \quad Azufre: no metal.

    \item a) Dalton ($\approx 1808$) $\rightarrow$ Thomson (1897) $\rightarrow$ Rutherford (1911) $\rightarrow$ Bohr (1913).
          \quad b) Rutherford descubrió que el átomo tiene un \textbf{núcleo} pequeño, denso y positivo, y que es mayormente espacio vacío: la mayoría de las partículas alfa atravesaban la lámina de oro sin desviarse, pero unas pocas rebotaban.
  \end{enumerate}

  \newpage

  \subsection*{Sección 4: Transformaciones y Ecuaciones Químicas}

  \begin{enumerate}[leftmargin=*, label=\textbf{\arabic*.}]
    \item a) F \quad b) Q \quad c) F \quad d) Q \quad e) F \quad f) Q

    \item a) Reactivos: metano (CH$_4$) y oxígeno (O$_2$). Productos: dióxido de carbono (CO$_2$) y agua (H$_2$O).
          \quad b) CH$_4$: gas; O$_2$: gas; CO$_2$: gas; H$_2$O: gas.
          \quad c) «Una molécula de metano gaseoso reacciona con dos moléculas de oxígeno gaseoso y produce una molécula de dióxido de carbono y dos moléculas de agua, ambas gaseosas.»

    \item a) $2\,\text{H}_2\text{O(l)} \xrightarrow{\text{electricidad}} 2\,\text{H}_2\text{(g)} + \text{O}_2\text{(g)}$.
          \quad b) Tres indicios: cambio de color permanente, producción de gas (burbujas) y emisión o absorción de calor/luz (también la formación de un precipitado o un cambio de olor).
  \end{enumerate}

\fi
\end{document}
